\documentclass[]{article}

\usepackage{mathtools}

\begin{document}
\section{Global}
\subsection{Objekte}
Energietypen: $i \in EC$ (Energy Carrier) Zeit: $t \in T$, Komponenten: $d \in D$
Insgesamt gibt es folgende Objekte:
\begin{itemize}
    \item Grid: $G \subset D$ externes Energienetz ("Slack") kann Energie Ein- oder Verkaufen
    \item Renewable Generation: $R \subset D$ lokale Erzeugung von Energie, wird immer eingespeist
    \item Demand: $D \subset D$ lokale Verbraucher mit bekannten Lastprofilen
    \item Converter $C \subset D$ Umwandlung von verschiedenen Energietypen z.B. Elektrolyse, BHKW
    \item Storage $S \subset D$ Speicherung von Energie (Batterie, Gasspeicher, \ldots)
\end{itemize}
\subsection{Komponenten}
Alle Komponenten: $p_{d,i,t}$ Leistung zum Zeitpunkt.
Vorzeichen +: Leistung geht aus der Komponente raus, -: Leistung geht in die Komponente rein.
Global eine Einheit für jeden Energietyp festlegen, bei Modellierung der Geräte beachten.
\subsection{Energiebalance}
\begin{equation}
    \sum_{D} p_{d,i,t} = 0 \quad \forall i \in EC \quad \forall t \in T
\end{equation}
\subsection{Objectives}
Summe der einzelnen Objectives der Komponenten. Eventuell noch global gewichten nach Wirtschaftlich/Co2/Fulfillment/\ldots
\section{(External) Grid}
Anschluss ans Netz, auch möglich für simpel modellierte eigene Erzeugung.
$G \subset D$ ist die Menge der Komponenten, die als Supply fungieren.
\subsection{Konstanten}
\begin{itemize}
    \item $c_{d,i,t}$: Preis per Energietyp, per Zeitpunkt im Ein-/Verkauf.
    \item $P_{d,i,t}^{+}$: Maximaler Einkauf per Energietyp, per Zeitpunkt
    \item $P_{d,i,t}^{-}$: Maximaler Verkauf per Energietyp, per Zeitpunkt (Negative Zahl)
\end{itemize}
\subsection{Constraints}
\begin{equation}
    p_{d,i,t} \leq P_{d,i,t}^{+} \quad \forall i \in EC \quad \forall t \in T
\quad \forall d \in S
\end{equation}
\begin{equation}
    p_{d,i,t} \geq P_{d,i,t}^{-} \quad \forall i \in EC \quad \forall t \in T
\quad \forall d \in S
\end{equation}
\subsection{Objectives}
Kosten für Energie/Einspeisung:
\begin{equation}
    \sum_{i=1}^{EC} \sum_{t=1}^{T} p_{d,i,t} \cdot c_{d,i,t}
\end{equation}

\section{Renewable Generation}
Lokale, Fixe Erzeugung von Energie. $G \subset D$
\subsection{Konstanten}
$P_{d,i,t}^{+}$: Einspeisung pro Energietyp, pro Zeitpunkt
\subsection{Constraints}
Einspeisung findet immer statt:
\begin{equation}
    p_{d,i,t} = P_{d,i,t}^{+} \quad \forall i \in EC \quad \forall t \in T \quad \forall d \in G
\end{equation}

\section{Demand}
Lokale Verbraucher mit fixen Lastprofilen. $D \subset S$
Wenn man die Bedarfsdeckung weglässt auch mit Renewable Generation zusammengefasst möglich.
\subsection{Konstanten}
$L_{d,i,t}$: Last per Energietyp, per Zeitpunkt (Negative Zahl)
\subsection{Constraints}
Bedarf wird gedeckt, mit Kostenfunktion unten weglassen.
\begin{equation}
    p_{d,i,t} = L_{d,i,t} \quad \forall i \in EC \quad \forall t \in T \quad \forall d \in D
\end{equation}


\section{Converter}
Umwandlung von verschiedenen Energietypen. $C \subset S$. Kann mehrere Inputs und Outputs haben.
\subsection{Konstanten}
\begin{itemize}
    \item $P_{d,i}^{max}$: Maximale Leistung per Energietyp - als Negative Zahl angeben, wenn Input
    \item $P_{d,i}^{min}$: Minimale Leistung per Energietyp - als Negative Zahl angeben, wenn Input
    \item $EC_{in}$: Menge der Energietypen, die als Input verwendet werden
    \item $EC_{out}$: Menge der Energietypen, die als Output verwendet werden
    \item $CF_{d,i}$: Umwandlungsfaktor per Energietyp (z.B. 10kWh Strom erzeugt $1m^3$ Gas: $E_{d,strom} = \frac{1}{10}$ und $E_{d,gas} = 1$) (Global festgesetzte Einheiten beachten)
    \item $R_{d}^{up}$: Maximale Änderung der Leistung pro Zeitschritt nach oben (ramp up) als Anteil der Maximalleistung
    \item $R_{d}^{down}$: Maximale Änderung der Leistung pro Zeitschritt nach unten (ramp down) als Anteil der Maximalleistung
\end{itemize}
\subsection{Variablen}
\begin{itemize}
    \item $sp_{d,t} \in [0,1]$ Setpoint als Anteil der maximalen Leistung zur Steuerung
    \item $a_{d,t}$ Binäre Entscheidungsvariable, ob der Prozess aktiv ist oder nicht
\end{itemize}
\subsection{Constraints}
Inputs kleiner Null:
\begin{equation}
    p_{d,i,t} \leq 0 \quad \forall i \in EC_{in} \quad \forall t \in T \quad \forall d \in C
\end{equation}
Outputs größer Null:
\begin{equation}
    p_{d,i,t} \geq 0 \quad \forall i \in EC_{out} \quad \forall t \in T \quad \forall d \in C
\end{equation}
Umwandlungsbilanz - hier ist Effizienz implizit enthalten:
\begin{equation}
    \sum_{EC_{in}} p_{d,i,t} \cdot CF_{d,i} - \sum_{EC_{out}}p_{d,i,t} \cdot CF_{d,i} = 0 \quad \forall d \in C \quad \forall t \in T
\end{equation}
Produziert Mindestenergie:
\begin{equation}
    P_{d,i}^{min} \cdot a_{d,t} \leq P_{d,i}^{max} \cdot sp_{d,t} \quad \forall i \in EC_{out} \quad \forall t \in T \quad \forall d \in C
\end{equation}
\begin{equation}
    P_{d,i}^{min} \cdot a_{d,t} >= P_{d,i}^{max} \cdot sp_{d,t} \quad \forall i \in EC_{in} \quad \forall t \in T \quad \forall d \in C
\end{equation}
Output:
\begin{equation}
    p_{d,i,t} = P_{d,i}^{max} \cdot sp_{d,t} \quad \forall i \in EC \quad \forall t \in T \quad \forall d \in C
\end{equation}
Setpoint und Aktivität sind parallel:
\begin{equation}
    0 \leq a_{d,t} - sp_{d,t} \quad \forall t \in T \quad \forall d \in G
\end{equation}
Vorher inaktiv: mindestens auf $P_{d,i}^{min}$ hochfahren, ansonsten Rampe beachten:
Hierfür: $sp_{min}$ (Nur zur Lesbarkeit, keine Optimierungsvariable)
\begin{equation}
    sp_{min} \cdot P_{d,i}^{max} = P_{d,i}^{min}
\end{equation}
\begin{equation}
    sp_{d,i,t} - sp_{d,i,t-1} \leq R_{d}^{up} \cdot a_{d,t - 1} \quad \forall i \in EC \quad \forall t \in T \quad \forall d \in G
\end{equation}
\begin{equation}
    sp_{d,i,t} \geq sp_{min} \cdot a_{d,t} \quad \forall i \in EC \quad \forall t \in T \quad \forall d \in G
\end{equation}
Ramp down:
\begin{equation}
    sp_{d,i,t-1} - sp_{d,i,t} \leq R_{d}^{down} \quad \forall i \in EC \quad \forall t \in T \quad \forall d \in G
\end{equation}
\section{Storage}
Gedacht als: Speicherung von Energie z.B. Batterie. $B \subset D$ Nur ein Energietyp als Input und Output.
\subsection{Konstanten}
\begin{itemize}
    \item $P_{d}^{+}$: Maximale Positive Leistung - Entladerate
    \item $P_{d}^{-}$: Maximale Negative Leistung - Laderate
    \item $\eta_{d} \in [0,1]$: Entladeeffizienz
    \item $\eta_{c} \in [0,1]$: Ladeeffizienz
    \item $C_{d}$: Kapazität
    \item $\Delta t$: Zeitschrittgröße
\end{itemize}
\subsection{Variablen}
\begin{itemize}
    \item $soc_{d,t} \in [0,1]$ State of Charge per Energietyp, per Zeitpunkt (Anteil der Kapazität)
    \item $sp_{d,t} \in [-1,1]$ Setpoint als Anteil der maximalen Leistung zur Steuerung des Ladens/Entladens ($<0$: Laden, $>0$: Entladen)
    \item $y_t \in \{0,1\}$ Hilfsvariable ob geladen oder entladen wird (ergibt sich aus $sp_{d,t}$) (1 = entladen, 0 = laden)
\end{itemize}
\subsection{Constraints}
Hilfsvariable $y_{t} = 1$ wenn entladen wird
\begin{equation}
    sp_{d,t} >= -(1 - y_{t}) \quad \forall t \in T \quad \forall d \in B
\end{equation}
$y_{t} = 0$ wenn geladen wird
\begin{equation}
    sp_{d,t} <= y_{t} \quad \forall t \in T \quad \forall d \in B
\end{equation}
Leistung kleiner als Maximale Entladerate
\begin{equation}
    p_{d,t} \leq P_{d}^{+} \quad \forall t \in T \quad \forall d \in B
\end{equation}
Leistung größer als Maximale Laderate
\begin{equation}
    p_{d,t} \geq - P_{d,t}^{-} \quad \forall t \in T \quad \forall d \in B
\end{equation}
Output kleiner als State of Charge
\begin{equation}
    y_{t} \cdot p_{d,t} \cdot (E_{discharge})^{-1} \cdot \Delta t\leq soc_{d,t} \cdot C_{d} \quad \forall t \in T \quad \forall d \in B
\end{equation}
Input kleiner als Kapazität und Maximaler Laderate
\begin{equation}
    (1 - y_{t}) \cdot p_{d,t} \cdot E_{charge} \cdot \Delta t \leq C_{d} \cdot soc_{d,t} \quad \forall t \in T \quad \forall d \in B
\end{equation}
Setpoints für laden und entladen:
\begin{equation}
    p_{d,t} = P_{d}^{+} \cdot sp_{d,t} \cdot y_{t} + P_{d}^{-} \cdot  sp_{d,t} \cdot (1 - y_{t})\quad \forall t \in T \quad \forall d \in B
\end{equation}
Hilfsvariable für absolute Ladung (nur zur Lesbarkeit)
\begin{equation}
    c_{d,t} = soc_{d,t} \cdot C_{d} \quad \forall t \in T \quad \forall d \in B
\end{equation}
Damit: Konsistenter State of Charge
\begin{equation}
    c_{d,t} = c_{d,t-1} - (p_{d,t} \cdot y_{t} \cdot \eta_{c} + (1 - y_{t}) \cdot p_{d,t} \cdot \eta_{d}^{-1}) \cdot \Delta t \quad \forall t \in T \quad \forall d \in B
\end{equation}
Wenn $t=0$ dann $p = 0$, initialen Wert von $soc$ nehmen.
\end{document}